% Nouveau EGA V (cote25): archive cote25, site121 (PDF page122).
% English translation of the checked French annotated master.
% AI-translated by GPT-6 Astra (OpenAI, ChatGPT Pro mode).
\scanpage{25/121}
\mspage{24}

will automatically be a \var{\deleted{Cartier}}{}{deletion}{ega25_25_121-v0001}\var{\inserted{divisor}}{divisor}{insertion}{ega25_25_121-v0002} relative to the base $\overset{\hand{\vee}}{P}$ , as
follows from the fact that U contains the \unsure{elements} associated with $\underline{O}_{X_s}$ , s $\hand{\in}$ S,
a condition stable under every base change and in particular
under S'=$\overset{\hand{\vee}}{P} \hand{\longrightarrow}$ S. Now it follows immediately from lemma \var{\inserted{4.2.}}{4.2.}{insertion}{ega25_25_121-v0003} \unsure{above} that D'
extends to a \var{\deleted{Cartier}}{}{deletion}{ega25_25_121-v0004} di\unsure{v}isor D \uline{iff} $\underline{N}'$ extends to an invertible
Module, or again $i'_{\hand{*}}(\underline{N}')$ is invertible. \var{\deleted{Therefore}}{}{deletion}{ega25_25_121-v0005} \var{\over{i}{I}}{I}{writing over}{ega25_25_121-v0006}t would be necessary to \unsure{write}
the end of the \unsure{proof} better in terms of nec and suff cond (without
going through it twice as I did), and also to \unsure{bring out} \var{\inserted{first}}{first}{insertion}{ega25_25_121-v0007} the
\var{\margin{$|$}}{\placedmargin{$|$}}{marginal note placed}{ega25_25_121-v0008}\var{\margin{\unsure{see also} \S 20.\illegible{1} and 20.11.}}{\placedmargin{\unsure{see also} \S 20.\illegible{1} and 20.11.}}{marginal note placed}{ega25_25_121-v0009}\uline{Corollary} \var{\inserted{3.4.}}{3.4.}{insertion}{ega25_25_121-v0010} \uline{to lemma}\hand{3.2}. Suppose g:X$\hand{\to}$S flat loc of finite presentation;
let Z be a closed subset of X such that x $\hand{\in}$ Z implies prof$_{\unsure{x|s}}$ $\underline{O}_{X_s} \hand{\geq} 2$
\var{\inserted[typed]{let U=X-Z and i:U$\hand{\to}$X be the can. immersion.}}{let U=X-Z and i:U$\hand{\to}$X be the can. immersion.}{insertion}{ega25_25_121-v0011}
(where s = g(x)), \var{\xout{Let $\underline{L}$ be a loc free Module on U = X-Z, then $\underline{L}'$}}{}{typewriter strike-out}{ega25_25_121-v0012}
\var{\xout{is isomorphic to the restriction}}{}{typewriter strike-out}{ega25_25_121-v0013} For every loc free Module \var{\inserted[typed]{of f.t}}{of f.t}{insertion}{ega25_25_121-v0014} $\underline{L}$ on X,
consider its restriction $\underline{L}' = \underline{L}\hand{|}U$ . Then \var{\inserted[typed]{a)}}{a)}{insertion}{ega25_25_121-v0015} the functor $\underline{L} \hand{\rightsquigarrow} \underline{L}'$ is
fully faithful, and for \var{\xout{some loc free Module $\underline{L}$}}{}{typewriter strike-out}{ega25_25_121-v0016} every $\underline{L}$, \var{\xout{one has}}{}{typewriter strike-out}{ega25_25_121-v0017}
the canonical hom $\underline{L} \hand{\to} i_{\hand{*}}(\underline{L}')$ is an isomorphism. \var{\inserted[typed]{For $\underline{L}$ to have rank n, it is n and s that $\underline{L}'$ does.}}{For $\underline{L}$ to have rank n, it is n and s that $\underline{L}'$ does.}{insertion}{ega25_25_121-v0018} b) Let $\underline{L}'$ be a
loc free Module on X; then $\underline{L}'$ is isomorphic to the restriction of a
loc free Module $\underline{L}$ iff $i_{\hand{*}}(\underline{L}')$ is loc free. c) Suppose $\underline{L}'$ is
the invertible Module associated with a \var{\deleted{Cartier}}{}{deletion}{ega25_25_121-v0019} divisor D'\hand{/}on U. Then
the condition considered in b) is also nec and sufficient for D'
to be the restriction of a \var{\deleted{Cartier}}{}{deletion}{ega25_25_121-v0020} divisor D on X, \var{\inserted[typed]{which will then be unique,}}{which will then be unique,}{insertion}{ega25_25_121-v0021}\var{\inserted{and equal to the schematic closure of D' in X.}}{ and equal to the schematic closure of D' in X.}{insertion}{ega25_25_121-v0022} For D' to be
a \var{\deleted{Cartier}}{}{deletion}{ega25_25_121-v0023} divisor relative to S, it is n. and s. \var{\inserted{that}}{that}{insertion}{ega25_25_121-v0024} D is one.

One simply uses the fact that every $\underline{L}$ satisfies the \var{\xout{cond}}{}{typewriter strike-out}{ega25_25_121-v0025}
hypotheses stated for $\underline{F}$ in lemma \hand{3.2}...

\uline{Corollary} \var{\inserted{3.5.}}{3.5.}{insertion}{ega25_25_121-v0026} Suppose \var{\inserted{the local rings of X factorial, (for example}}{the local rings of X factorial, (for example}{insertion}{ega25_25_121-v0027} X regular\hand{)} \var{\deleted{(for example S regular, and X smooth over
S)}}{}{deletion}{ega25_25_121-v0028}; Then the map $\hand{\unsure{\phi|\emptyset}} \hand{\to} \hand{\unsure{\psi|\emptyset}}$ is bijective. In particular, if X is
a regular prescheme loc of finite type over a field k, and P a projective
bundle over k, there is a bijection between the set $\hand{\unsure{\phi|\emptyset}}$ of
\tnote{25-121-1}{The doubtful letters in French ``rédiger'', ``démonstration'', and ``dégager'' are represented by doubts on ``write'', ``proof'', and ``bring out''. ``Iff'' renders ``sss'', si et seulement si. ``Nec and suff cond'' and ``n and s'' retain the abbreviated necessary-and-sufficient wording. ``Loc'', ``can.'', ``f.t'', and ``hom'' mean locally, canonical, finite type, and homomorphism; f.t translates French t.f. The canceled ``con'' remains a fragment as ``cond''. The reference4.2 and the X in part(b) are preserved.}

\var{\margin{121}}{\placedmargin{121}}{marginal note placed}{ega25_25_121-v0029}
